
\chapter{Conclusion}
\label{chap:conclusion}

This thesis examined audited semantic propagation in an LLM-agent society. It
connected a measurement instrument to persona composition, topology, a separate
complex network, and Pfeffer's firestorm account. The evidence supports a
bounded mechanism study. It does not support a model of real Twitter users.

\section{Answers to the research questions}

\paragraph{RQ1: measurement.}
The continuous and dual-discrete auditors produced different levels and weakly
aligned topology rankings. Dual-discrete MI exceeded continuous MI on all eight
topology aggregates in both identity arms. The mean topology gaps were 0.8201
for H and 1.0450 for He. These values show instrument dependence. They do not
show that one auditor measured the true amount of misinformation. Continuous
and discrete MPR must therefore remain separate.

\paragraph{RQ2: persona composition.}
Explicit composition described the results more consistently than the H/He
label. Homogeneous conspiracy-family means were 2.2194 continuous and 3.1691
dual-discrete. Science and environment means were 0.1118 and 0.8687. The
conspiracy-heavy \texttt{mix\_00} also exceeded conspiracy-free
\texttt{mix\_02} under both instruments. These contrasts apply to the executed
prompts and model. They do not establish a general effect of human diversity.

\paragraph{RQ3: topology and tipping.}
Topology conditioned the composition contrast, but no universal topology
ranking emerged. Continuous He--H differences changed sign across the eight
graphs. Dual-discrete differences were positive on all eight. Dead-cell rates
also varied by topology, and live-only means describe surviving cascades.
The observed \(k^{*}\) values are eight-tick, instrument-specific outcomes.
They are not empirical ignition times.

\paragraph{RQ4: complex-network generalisation.}
The 63-node D-net preserved an empirically motivated hashtag structure and
supported a controlled composition comparison. Conspiracy-only H exceeded the
mixed assignment for both articles under both instruments. Structural
correspondence remained partial. Breadth met the selected tolerance, whereas
depth and structural virality did not. The Digital Twin Fidelity Score was
0.2754 and validation failed. D-net is therefore a separate boundary test, not
a validated retweet cascade or a source of empirical Twitter MPR.

\section{Contributions and boundary}

The thesis contributes a transparent theory-to-design ledger, a 1,728-cell
climate factorial, an explicit measurement-sensitivity analysis, a
composition-conditioned interpretation, and an honest 63-node boundary test.
The canonical harvest contains 288 configurations, 1,728 cells, and 292
hatched-dead cells. It uses \texttt{gpt-4o-mini}, eight hops, and one run per
configuration. The C5 ``combined'' result is a table union. It is not a
physical supergraph.

Pfeffer, Zorbach, and Carley identified seven Outlook factors
\parencite{pfeffer2014firestorms}. The six-knob vocabulary in this thesis is a
computational remapping. Cross-media dynamics remained held. Empirical timing
was unavailable. The engine's ternary actions did not implement the original
binary-choice factor. These omissions prevent a seven-factor confirmation.

\section{Outlook}

The next study should prioritise validation over scale. Independent graph and
model seeds should support uncertainty estimates. Matched persona compositions
should isolate diversity from prevalence. Independent human coders should
assess factuality and realism under an approved protocol. Longer horizons
should test censoring at tick eight. Lawfully obtained timestamped interaction
data would permit a genuine temporal and cascade comparison. Until those
conditions are met, this testbed is best used to probe mechanisms and expose
measurement assumptions.
